\subsection{Overview}

Manure and other organic amendments sit at the intersection of two functions of subsurface (tile) drainage systems: they enrich soils with nutrients and organic matter, and they introduce fecal microorganisms, resistance genes, hormones, and trace elements that drainage networks can rapidly route to receiving waters. The evidence assembled here establishes several consolidated conclusions. First, the \emph{method} of manure application and the \emph{moisture state} of the soil, rather than manure presence alone, largely govern how much fecal contamination reaches tile drains, with preferential (macropore) flow acting as the dominant transport mechanism \parencite{Truhlar2015,Thiagarajan2007,Collins2007,Ding2021}. Second, manure and organic amendments deliver measurable agronomic and soil-structural benefits---reclamation of saline--sodic soils, higher yields, and, when combined with subsurface drainage and biochar, improved soil physical properties and nitrogen-use efficiency---though some physical benefits are transient \parencite{Qadir2001,Tian2023b,Liu2025f,Xiong2025,Brikmans2025}. Third, drainage systems and their engineered analogues (denitrifying bioreactors, constructed wetlands, reservoirs, and reconstructed oxbows) host distinct and dynamic microbial communities and can substantially attenuate nitrate and pathogen loads \parencite{Andrus2014,Porter2015,Hathaway2015,Knorr2023,Schilling2017,Haverstock2017}. Fourth, drainage water can concentrate and mobilise trace elements---most infamously selenium at Kesterson, but also arsenic, boron, uranium, and micronutrients---with documented consequences for wildlife \parencite{SAIKI1987,Ohlendorf2002,Ryu2010,Gao2007,Tayfur2010,Gul2013}. Finally, subsurface drainage reshapes the ecology of agricultural landscapes, from soil fauna to birds, amphibians, and semi-aquatic species \parencite{Nuutinen2001,Marja2013,Dreyfuss2018,Campbell2023}. This section synthesises the mechanisms and quantified outcomes underpinning each of these lines of evidence.

\subsection{Manure Application Method, Soil Moisture, and Microbial Transport to Tile Drains}

\subsubsection{Application method controls the transport of fecal bacteria}
Controlled comparisons demonstrate that how manure is placed on the field determines the microbial signal reaching drains. \textcite{Truhlar2015} found that roughly 45\% of \textit{E.\ coli} isolated from manure produced curli (an adhesive amyloid fibre) before field application, and that curli-producing \textit{E.\ coli} were far more abundant in tile drains from farms that surface-applied manure than in the drain of a farm that incorporated it. The same pattern held in soil-column leachate: the proportion of curli-producing \textit{E.\ coli} was greater under surface application than under incorporation or no manure, while total \textit{E.\ coli} recovered from leachate was higher in both manure treatments than in controls but did not differ between surface-applied and incorporated manure. Because the column manure came from a single source, and because the functional gene \textit{csgA} was present in all seven tested strains regardless of curli phenotype, the authors ruled out strain composition as the sole explanation and concluded that application method itself shapes the phenotype of \textit{E.\ coli} transported to drains---a result consistent across two surface-applied periods with different rainfall \parencite{Truhlar2015}. Complementary indicator work found that faecal coliforms, coliphages, and \textit{Salmonella} behaved differently: \textit{Salmonella anatum} was not detected in any of 36 leachate samples, male-specific coliphages were detected in all 36 samples, somatic coliphages were generally lower, and phage counts exceeded bacterial indicator counts in all plots; indicator numbers were similar in manure-treated and control plots, peaking in weeks 1--2 and declining thereafter \parencite{Malik2004}.

\subsubsection{Tillage, preferential flow, and pathogen loads}
Tillage governs pathogen delivery mainly through its effect on macroporosity and flow partitioning. Across two Canadian sites, zero tillage (ZT) generally produced higher subsurface flow volumes and greater \textit{E.\ coli} loads than conventional tillage, which the authors attribute to macropore development and preferential flow \parencite{Thiagarajan2007}. At Streets Ridge, combined-drainage \textit{E.\ coli} loads showed no overall significant tillage effect and the direction of differences switched between years, but when surface and subsurface drainage were separated, subsurface drainage under ZT contributed the greatest \textit{E.\ coli} losses in 2003--2004. At BEEC, tillage did not significantly affect \textit{E.\ coli} concentrations, yet subsurface-drainage loads were significantly higher under ZT in 2003--2004. Most samples exceeded drinking-water guidelines, and recreational-guideline exceedances were concentrated soon after manure application and were greater at Streets Ridge than at BEEC \parencite{Thiagarajan2007}. Hydrologic partitioning by tillage was also evident for manure-borne hormones: \textcite{Jenkins2009} reported that no-till plots had greater drainage and lower runoff than conventional-till plots in June 2002, and that November 2001 runoff loads of estradiol and testosterone were lower from no-till than conventional-till plots.

\subsubsection{Soil moisture, water table, and drain-flow contamination}
Antecedent soil-moisture state and water-table position exert first-order control on contamination episodes. In New Zealand pastoral systems, effluent applied to wet soil produced drain-flow \textit{E.\ coli} of $10^{5}$--$10^{7}$ per 100\,ml, whereas a greater soil-moisture deficit reduced concentrations to about $10^{4}$ per 100\,ml \parencite{Collins2007}. Peak drain-flow \textit{Campylobacter} exceeded $10^{3}$ per 100\,ml following effluent application at negligible soil-moisture deficit, and where preferential flow occurred, drainage concentrations approached those of the applied effluent. Over reported drainage events, \textit{E.\ coli} yield from a 1080\,m$^{2}$ plot ranged between $10^{8}$ and $5\times10^{10}$; installing subsurface drains on Pallic Soils reduced surface runoff by 60--80\%, redirecting water (and its microbial load) to the subsurface pathway \parencite{Collins2007}. Winter reuse of undisinfected wastewater in Ohio showed the same water-table dependence: during a non-irrigated control winter, \textit{E.\ coli} was detected in 2 of 6 subsurface samples when the water table was within the drain depth and briefly after it dropped below the drains; during the winter irrigated season, \textit{E.\ coli} was detected three times, always when the water table rose above the drain under cold, wet soil conditions \parencite{Ding2021}. Soil moisture and temperature were significantly positively related at a given depth, sunlight exposure of wastewater in metal pans for 2--3\,h produced up to 3-log \textit{E.\ coli} reductions, and \textit{E.\ coli} was not detected in groundwater 15\,m from the field---leading the authors to conclude that winter wastewater reuse poses a public-health risk to surface water via subsurface drainage, with irrigation rate and disinfection as key management levers \parencite{Ding2021}. At the pathway scale, \textcite{Collins2007} further documented that direct herd crossings of the Sherry River temporarily raised \textit{E.\ coli} to more than 100 times background (about $10^{11}$ \textit{E.\ coli} deposited during one crossing over 25 defecation events), and estimated a 60--75\% probability of \textit{Campylobacter} infection over an irrigation season from observed groundwater concentrations.

\subsubsection{Antibiotic resistance genes and hormones carried by manure}
Manure is a reservoir of antibiotic resistance genes (ARGs) that persist through drained soils into drainage water. Shotgun metagenomics identified 320 unique ARGs, with 173 in manure, 176 in soil, and 85 shared; manure treatment significantly shifted the soil and drainage-water resistomes away from the control-soil resistome \parencite{Smith2019b}. Of 88 ARGs unique to the manure resistome, 75 appeared in manure-treated soil and 69 in manure-treated effluent, and 25 showed strong persistence across at least half of samples and all time points. Manure-derived ARGs dominated the abundance in manure-treated effluent (manure-enriched 0.612, water-enriched 0.308, manure persisters 0.071), with \textit{tet}(33) the most abundant persisting sequence, rising over time to become most abundant by day 108. Whereas untreated-soil ARGs were dominated by multidrug-efflux genes, many manure-derived ARGs were associated with mobile genetic elements in food- and water-borne pathogens---evidence that manure-derived genes have strong potential to persist in soil and water in tile-drained landscapes \parencite{Smith2019b}. For steroid hormones, poultry litter did not increase the flow-weighted concentration or load of estradiol or testosterone in subsurface drainage or surface runoff relative to conventional fertiliser/control treatments; both hormones were detected in drainage and runoff but total loads were generally not different among treatments, soil concentrations declined with depth (indicating stratification), and mean bulk-soil testosterone was greater in poultry-litter-amended plots while soil estradiol did not differ \parencite{Jenkins2009}.

\subsection{Organic Amendments and Soil Amelioration in Drained Soils}

\subsubsection{Manure and organic fertiliser for saline--sodic reclamation}
Farm manure is an effective amendment for reclaiming saline--sodic soils. In Pakistan, active desalinisation and desodication occurred in all treatments with non-significant differences among them; after the final wheat crop the 1.2\,m profile reached ECe $\approx 7\pm0.5$\,dS\,m$^{-1}$ and SAR $\approx 15\pm2$, indicating comparable amelioration by phytoremediation, gypsum, sulphuric acid, and farm manure \parencite{Qadir2001}. Farm manure nonetheless produced the best yields, giving 4188\,kg\,ha$^{-1}$ for the first wheat crop versus 3462\,kg\,ha$^{-1}$ (100\% gypsum), 2854\,kg\,ha$^{-1}$ (phytoremediation), 2786\,kg\,ha$^{-1}$ (sulphuric acid), and 2754\,kg\,ha$^{-1}$ (50\% gypsum); infiltration rate declined by about 18\% (from $2.03\times10^{-6}$ to $1.67\times10^{-6}$\,m\,s$^{-1}$) with non-significant treatment differences, and phytoremediation required minimal capital input \parencite{Qadir2001}. Under monsoonal Russian conditions, drainage and organic fertiliser had stronger effects on soil pore volume than biochar, with organic fertiliser increasing porosity in 2019 \parencite{Brikmans2025}. A North Dakota case study showed that a legacy of combined subsurface drainage, no-till, gypsum, and cover crops protected crops from salinisation through soil processes not revealed by soil indices alone: in the soil-health-legacy zone, 6 of 9 chemical, 5 of 14 microbial, and none of 3 physical properties differed between areas with dead/damaged plants and living plants, with higher ECe, Olsen-P, Ca, Na, and SAR (and lower pH) where plants were dying, and lower protein, POXC, arbuscular-mycorrhizal fungi, and G+:G- ratio \parencite{Daigh2025}.

\subsubsection{Biochar combined with subsurface drainage on saline--alkali soils}
A consistent body of Chinese coastal-soil experiments shows that pairing biochar with subsurface drainage improves saline--alkali soils beyond either practice alone. Subsurface drainage plus biochar (B-S) significantly reduced soil electrical conductivity and water content relative to no drainage, and increased soil total nitrogen, available phosphorus, nitrate, and organic matter relative to both control and drainage-alone; B-S raised alkaline phosphatase and urease activity, increased bacterial alpha diversity (Chao1, Shannon), and shifted community composition (Proteobacteria dominant), with available phosphorus and organic matter identified as dominant drivers of enzyme activity and bacterial abundance \parencite{Tian2023b}. Narrower drain spacing and higher biochar rates lowered bulk density and raised porosity, field capacity, and soil moisture in the plow layer, affected the entire 0--80\,cm profile, and reduced salinity---with drainage a stronger salinity control than biochar; the S1C3 combination (6\,m spacing, 15\,t\,ha$^{-1}$ biochar) gave the most satisfactory salt reduction and, with the S1 spacing, the largest yield gains \parencite{Liu2025f}. Using $^{15}$N-labelled urea, \textcite{Xiong2025} showed that narrower spacing and more biochar promoted $^{15}$N absorption across all three alfalfa harvests (drainage having the more persistent effect); S1C3 produced the highest uptake (6.8, 15.0, 12.8\,kg\,ha$^{-1}$) and raised $^{15}$N use efficiency to 46.0\% (versus 21.6--28.4\% in controls), with 21.6--46.0\% of applied $^{15}$N taken up by crops, 38.6--67.5\% retained in soil, and 8.5--18.1\% lost. In contrast, biochar's physical benefits in clay-loam soils were transient: dose-dependent decreases in bulk and aggregate density in 2018 dissipated by 2019, total porosity showed no sustained biochar effect, and intra-aggregate pore structure (persistent mesopore dominance) resisted amendment, consistent with biochar occupying inter-aggregate spaces temporarily rather than reshaping intra-aggregate pores \parencite{Brikmans2025}.

\subsection{Soil Physical and Hydraulic Properties under Drainage and Management}

Management practices exert clearer effects on soil structure than tile drainage does in several settings. In South Dakota silty clay loam, no-till significantly increased bulk density, soil organic matter (SOM), and wet aggregate stability (WAS) while decreasing field-saturated hydraulic conductivity (Kfs) relative to conventional tillage; more diverse rotations decreased bulk density and increased SOM and WAS; and cover crops and tile drainage had minimal or no significant impact on SOM, bulk density, WAS, or Kfs \parencite{Saha2024}. The tillage$\times$rotation interaction was significant (e.g., higher Kfs under conventional tillage than no-till in the 2-year rotation), bulk density was negatively correlated with SOM and WAS, SOM was strongly negatively correlated with pH, WAS was positively correlated with SOM, and stacking no-till with cover crops had no significant added effect \parencite{Saha2024}. Grass ley consistently improved surface aggregate stability of both drained and undrained soils relative to spring barley underseeded with clover, with improvement significantly correlated with rising soil organic carbon; the grass/drained treatment gave the greatest structural improvement, while the grass/undrained treatment matched it only when the soil was dry, and aggregate stability varied seasonally (lower in early spring, higher in fall) with soil water content \parencite{Hermawan1996}. On an Australian sodic soil, a drainage system functioned adequately (rapid post-rainfall drainage, clear drainage water), sunflower seed yield averaged 1.2\,t\,ha$^{-1}$ and was unaffected by drainage or gypsum, but wheat establishment was much lower on the undrained fallow, no-gypsum treatment (45 versus 214 plants\,m$^{-2}$); gypsum raised EC by $\sim$0.2\,dS\,m$^{-1}$, lowered pH from 5.8 to 5.5, minimised dispersed clay, and air-filled porosity was lowest on fallow, no-gypsum plots (8.3\% versus 15.1\% after sunflowers) \parencite{GARDNER1991}. Drainage itself altered texture and pore structure in a Russian silt loam, increasing the silt fraction, reducing the sand fraction, and driving opposite seasonal shifts in pore volume \parencite{Brikmans2025}.

\subsection{Perennial Crops, Rotations, and Nutrient/Water Loads in Tile Drainage}

Introducing perennials into annual rotations reduces subsurface nutrient export, with the strongest benefits during high-loading periods. A before--after control--impact (BACI) analysis in Ohio found that alfalfa significantly affected mean monthly subsurface discharge at Site A ($F=4.3$, $p=.04$) but not Site B; alfalfa significantly reduced subsurface NO$_3$-N, TN, and DRP loads at Site A but not Site B, significantly reduced TP at both sites, and reduced the variability of discharge and nutrient loads (smaller post-planting IQR increases in impact fields) \parencite{Arrueta2022}. Seasonally, alfalfa lowered median weekly NO$_3$-N loads in spring and winter at Site A while showing no significant seasonal effect at Site B \parencite{Arrueta2022}. A storm-event analysis at the same sites---grouping events into dry, wet, and large types---found alfalfa had little effect on discharge but significantly lowered subsurface NO$_3$-N and TN across all event types at Site A; during large events, DRP and TP loads were lower in the alfalfa field at \emph{both} sites, identifying large storms as the most important period for phosphorus-export reduction \parencite{Arrueta2025}. Collectively these results support partial, targeted conversion of marginal cropland to perennial cover to cut nutrient loads (especially dissolved and particulate P) and increase system resilience, ideally alongside reduced fertiliser rates and gypsum \parencite{Arrueta2022,Arrueta2025}. Drainage adequacy is itself indexed by perennial performance: on a sloping Cabot silt loam, subsurface drains (especially at 30.5\,m spacing) improved alfalfa growth and survival, percent alfalfa in summer cuttings correlated with plant density ($r=0.93$), percent winterkill correlated with spring ($r=0.79$) and fall ($r=0.84$) soil water, and draining to the 0.1-bar water content rendered winterkill negligible ($\approx$10\%), such that 30.5\,m spacing would ensure $\geq$50\% survival under severe conditions \parencite{BENOIT1967}. At the catchment scale, crop composition modulated phosphorus loss: summer crops occupied 63\% of the agricultural area and winter crops 34\%, summer-rape area correlated positively with summer P loss, winter-wheat area correlated strongly negatively with summer P load, grassland area correlated with winter P load, and no significant correlation linked N load to crop area \parencite{DambenieceMigliniece2018}. Long-term fertilisation and liming further shaped productivity: unlimed soil acidity changed with fertiliser rate (from $+0.3$ pH units unfertilised to $-1.1$ units at the highest rate over 34 years), liming raised pH to about 6.0--6.4, and fertiliser rate was the main control on perennial-grass productivity and sward botanical composition \parencite{Vigovskis2016}.

\subsection{Microbial Communities in Drainage Systems and Denitrifying Bioreactors}

Microbial communities in drained soils, bioreactors, and wetlands are spatially structured, temporally dynamic, and strongly shaped by the contemporary environment. In agricultural soil, the general bacterial community showed significant spatial structure by depth and transect, whereas the denitrifying community did not; directional geostatistics gave bacterial correlation ranges of 5.6\,m overall, 10.4\,m along the flow path, and 10.7\,m at 1.52\,m depth, supporting efficient sampling design (formalised in the \texttt{COMMUNITYCORRELOGRAM} R package) \parencite{Andrus2014}. In denitrifying bioreactors, depth and season shaped bacterial, fungal, and denitrifying communities; communities at 1.52\,m were more temporally variable than at 0.76\,m; bacterial communities cycled with a $\sim$125-day correlation time and roughly annual periodicity; denitrifiers recurred at $\sim$300 days; and fungi were more stable ($\sim$150--200-day autocorrelation, no significant annual cycling), with seasonal flow/moisture (Jan--Jun) and temperature/precipitation (Jul--Dec) gradients evident \parencite{Porter2015}. Comparing habitats, bacterial and clade-I \textit{nosZ} denitrifier communities were distinct among soils, bioreactors, and wetlands (PERMANOVA: significant location and habitat effects), with bioreactor communities more similar to wetlands than to fields---evidence that favourable contemporary environmental conditions can be engineered to develop effective denitrifier assemblages \parencite{Hathaway2015}. In a Canadian agricultural stream network, source locations (tile drains and the drainage ditch) had the highest N and P cycling functional potential and the most abundant nitrification/denitrification genes, storm sampling was associated with elevated nutrient concentrations and functional capacity, archaeal N-cycling capacity and NH$_4$-N rose after harvest, microbial P-mobilisation potential tracked P availability, and water passing through the wetland generally showed reduced nutrient concentrations and functional capacity, though the wetland response was variable \parencite{Knorr2023}.

\subsection{Drainage Water Treatment: Wetlands, Reservoirs, and Reconstructed Oxbows}

Engineered receiving features can substantially attenuate drainage-borne nitrate, though their effectiveness for pathogens is more variable. A reconstructed Iowa floodplain oxbow fed primarily by tile drainage retained nitrate with May--September efficiencies of 44\% (2014) and 47\% (2015) by N:Cl ratio, with retention greatest in late summer/early fall (monthly maxima of 92\% in August and 78\% in September 2015); nitrate concentrations were highest in tile water, similar in upgradient groundwater and the oxbow, and lowest downgradient, and groundwater contributed only a small fraction of discharge---placing oxbow performance on par with bioreactors, wetlands, and saturated buffers as a viable tile-drainage BMP \parencite{Schilling2017}. A Canadian water-and-reservoir-for-irrigation (WRI) constructed wetland reduced mass loads by 67.6\% for nitrate-nitrogen and 63.3\% for \textit{E.\ coli}; outflow nitrate averaged 1.2\,mg\,L$^{-1}$ (below the drinking-water guideline but exceeding the aquatic-life guideline on 25\,d), outflow \textit{E.\ coli} averaged 178\,CFU\,100\,mL$^{-1}$ (exceeding the irrigation guideline on day 119, higher in the warm season), measured residence time was 15.0\,d (design 14.5\,d), and areal uptake rates were 8.0\,m\,y$^{-1}$ (nitrate-N) and 6.4\,m\,y$^{-1}$ (\textit{E.\ coli}); the wetland captured 15{,}600\,m$^{3}$ (14{,}064\,m$^{3}$ after evapotranspiration, $\approx$780\,mm over 1.8\,ha), and a relatively small wetland (5.7\% of drainage area) sufficed \parencite{Haverstock2017}. Not all wetlands attenuate pathogens, however: a New Zealand wetland was consistently a net exporter of \textit{E.\ coli}, with loads increasing by a median 10.2-fold (range 1.5--26$\times$) across ten flow events (median inflow $3.8\times10^{8}$ versus outflow $2.5\times10^{9}$\,MPN/event), and monthly monitoring likewise showed higher outflow than inflow concentrations \parencite{Stott2023}. Selenium-removal wetlands treating drainage water in California retained about 59\% of inflow Se (outflow 35\%, seepage 4\%, volatilisation 2\%), achieving 21--55\% concentration reduction and 48--76\% mass removal in 2000, with the nonvegetated cell least effective---indicating that vegetation-promoted reducing conditions and retention in the organic detrital layer and surface sediment drive removal \parencite{Gao2003}.

\subsection{Trace Elements and Salinity in Agricultural Drainage Water}

\subsubsection{Selenium and the Kesterson legacy}
The Kesterson Reservoir remains the definitive case of trace-element hazard from agricultural drainage. Waterborne selenium at Kesterson averaged 0.122\,mg\,L$^{-1}$---about 230 times the $<$0.001\,mg\,L$^{-1}$ reference at Volta---and mean tissue selenium ran 10--200 times higher, biomagnifying from water to sediment to plants to animals; detritus reached 440\,$\mu$g\,g$^{-1}$ dry weight and some mosquitofish 370\,$\mu$g\,g$^{-1}$, levels judged probably toxic to sensitive consumers and consistent with the observed dead and deformed bird embryos \parencite{SAIKI1987}. \textcite{Ohlendorf2002} documented that 40\% of 347 nests in 1983 held one or more dead embryos and 20\% held severely deformed embryos or chicks; across 1983--1985 at least 39\% of 578 nests contained a dead or deformed embryo; bird-liver selenium was $\sim$10$\times$ and egg selenium 20--30$\times$ the Volta reference; food-chain organisms averaged 20--332\,$\mu$g\,g$^{-1}$ dry weight ($\sim$100$\times$ Volta); inflow selenium averaged 300\,$\mu$g\,L$^{-1}$; and coots showed mortality, reproductive failure, and selenium-toxicosis lesions. These integrated field and laboratory studies fulfilled the equivalent of Koch's postulates, identified selenium as the sole causative agent, and are regarded as a gold standard of retrospective ecological risk assessment \parencite{Ohlendorf2002}. Remediation drew on natural attenuation: groundwater beneath Kesterson was protected by biogeochemical immobilisation (selenium exceeded 10\,$\mu$g\,L$^{-1}$ in only 16 wells in 1985, clustered in two areas, within a $\sim$650\,ha drainage plume fed by $\sim$430\,ha-m\,yr$^{-1}$ of seepage during 1981--1986), and a risk-based strategy left the selenium inventory in place and filled low-lying areas to prevent wildlife-attracting ephemeral pools \parencite{BENSON1993}.

\subsubsection{Arsenic and uranium in evaporation ponds}
In agricultural evaporation ponds, evapoconcentration and redox chemistry jointly control trace-element mobility. Total dissolved arsenic rose from 108--132\,$\mu$g\,L$^{-1}$ at the central inlet to as high as 2{,}099\,$\mu$g\,L$^{-1}$ toward terminal cells; As increased with electrical conductivity yet arsenate As(V) remained dominant ($>$67\% of dissolved As in all cells, $>$95\% at the inlet), while reducing conditions toward terminal cells and in deeper layers lowered dissolved oxygen and nitrate and raised As(III), Fe(II), and sulfide \parencite{Ryu2010}. Sequential extraction showed carbonate-associated As as the largest sediment sink under both oxidising and reducing conditions, with soluble and adsorbed fractions much larger in the reduced terminal cell, so arsenic mobility was controlled by evaporation and by high-DOC reducing conditions \parencite{Ryu2010}. \textcite{Gao2007} corroborated this: pond-water As increased linearly with chloride (behaving more conservatively than selenium), speciation shifted toward As(III) and organic/methylated forms along the flow path and in reducing cells, sediments accumulated As but not strongly enough to prevent rising dissolved concentrations, and DOC increased with chloride and reducing conditions. Uranium showed analogous behaviour: sediment U decreased abruptly with depth (concentrated at the surface), was greatest in the most saline ponds, and correlated with sediment total organic carbon, implying association with algal organic matter; the measured Pond-14 K$_d$ of $\sim$50 far exceeded the isotherm-predicted 2, salt crusts held up to 5.8\,mg\,U\,kg$^{-1}$, and 75\% of surface-sediment U was oxidised versus 25\% reduced \parencite{Duff1997}.

\subsubsection{Boron and micronutrient balances}
Boron transport in agroforestry soils was successfully simulated by a two-dimensional finite-element model (mean absolute error 1.5\,mg\,L$^{-1}$, relative error 6.5\%), which reproduced Mendota field data for 1990 and 1992 and captured the effects of Langmuir constants: higher adsorption equilibrium constant, maximum adsorptive capacity, or root-absorption coefficient lowered soil-solution boron, while higher irrigation-water boron raised it, and adsorption was greater in silt loam than in sandy or clay loam \parencite{Tayfur2010}. Boron isotopes offer a manure/fertiliser tracer: hog manure held 0.4--8.2\,mg\,L$^{-1}$ B (mean 2.9) with $\delta^{11}$B of 7.2--11.2\textperthousand, cattle-feedlot runoff averaged 23.2\textperthousand, and Cobb River subsurface drainage averaged 0.06\,mg\,L$^{-1}$ B with $\delta^{11}$B of 5.3--15.1\textperthousand; the comparatively narrow hog-manure range makes $\delta^{11}$B useful for tracing feedlot contamination and mixing between water sources \parencite{Komor1997}. Micronutrient budgets can run negative under intensive drainage: in Pakistan, drainage-water removal of Cu, Fe, Mn, and Zn exceeded inputs from irrigation and fertiliser at both sites, giving negative balances and net soil depletion; Cu and Mn losses exceeded Fe and Zn (Zn least), losses were generally greater in maize (Kharif) than wheat (Rabi), and extractable concentrations rose with depth \parencite{Gul2013}.

\subsubsection{Major ions and salinity}
The ionic composition of central-Iowa drainage water was dominated (in decreasing order) by bicarbonate, calcium, magnesium, nitrate, chloride, sodium, and sulfate \parencite{Zimmerman2017b}. Bicarbonate, calcium, and magnesium were higher post-growing season while nitrate and chloride were higher in-season; corn-based systems had greater magnesium and chloride and lower calcium than soybean; bicarbonate correlated positively with calcium and negatively with nitrate and chloride; and flow, pH, and temperature did not differ among cropping systems, so they did not explain the ionic differences \parencite{Zimmerman2017b}. Salinity itself drove nonlinear soil responses in a North Dakota case study, where NO$_3$-N increased nonlinearly and Olsen-P linearly with salinity, while POXC, protein, and arbuscular-mycorrhizal fungi decreased \parencite{Daigh2025}; and subsurface drainage proved a stronger salinity control than biochar across the 0--80\,cm profile in coastal Chinese soils \parencite{Liu2025f}.

\subsection{Ecological Responses to Subsurface Drainage}

\subsubsection{Soil fauna}
Subsurface drainage patterns measurably shape earthworm populations. \textit{Lumbricus terrestris} abundance was typically higher above drains than between them---above the pipes the median individual count was twice as high and total fresh mass five times as high, with mean differences of 2.5 individuals\,m$^{-2}$ (95\% CI 1.0--4.0) and 6.6\,g\,m$^{-2}$ (95\% CI 3.6--9.6)---and stained burrows reached the 1\,m drain depth and below, so drainage pattern is a factor affecting population density and should inform sampling design (with worm activity potentially affecting subdrain function) \parencite{Nuutinen2001}. A long-term inoculation study confirmed both the above-drain preference (significantly higher densities above tile lines under no-till and ploughing) and strong management and dispersal gradients: by 2009 density reached 41 individuals\,m$^{-2}$ near margins versus 13 inside constantly unploughed plots, plot margins exceeded no-till which exceeded ploughing, density and fresh weight declined with distance from the inoculation area (14 to 1 individuals\,m$^{-2}$ from 5--9\,m to 56--60\,m), and an indicative colonisation rate of 4.6\,m\,yr$^{-1}$ was estimated---underscoring that undisturbed margin habitat and long-term ($>$10-year) monitoring are decisive for evaluating inoculation success \parencite{Nuutinen2011}.

\subsubsection{Birds, amphibians, and semi-aquatic species}
Drainage type is associated with farmland-bird outcomes. In this first long-term, large-scale Finnish comparison, skylark populations increased significantly on open-drained fields (ODF) but were stable on subsurface-drained fields (SDF); meadow pipit increased significantly on ODF (its SDF decline non-significant) and bred in significantly greater numbers on ODF; both species favoured ODF over SDF, with abundance not differing statistically between drainage types overall but territories fewer on arable than grassland fields \parencite{Marja2013}. Tile-drainage-associated pesticides posed hazards to amphibians: tile wetlands had cumulative pesticide concentrations 4.7 times higher than reference wetlands (neonicotinoids dominant), with clothianidin $\sim$8$\times$ and thiamethoxam nearly 9$\times$ higher (thiamethoxam detected at all tile wetlands and no reference wetlands); juvenile frogs from tile wetlands carried $\sim$2.5$\times$ more imidacloprid in brain tissue and nearly 2$\times$ more desnitro-imidacloprid in liver, and reference-wetland frogs had slightly longer standardised cerebellum lengths---the first field quantification of imidacloprid in juvenile amphibian brains and evidence that these chemicals can alter brain structure \parencite{Campbell2023}. Subsurface drainage of meadows was also identified, alongside ditch cleaning and gyro-crushing, as a leading cause of decline in the semi-aquatic snails \textit{Galba truncatula} and \textit{Omphiscola glabra}, whose occupied habitats fell by 30\% and 38\% respectively by 2016--2017, with the largest losses in vernal pools and significant reductions in drainage-furrow habitat area \parencite{Dreyfuss2018}.

\subsubsection{Vegetation and restored aquatic habitats}
Over 24 years, drained UK grassland shifted from a \textit{Lolium perenne}-dominated community with \textit{Juncus} patches toward one of \textit{Lolium perenne} and \textit{Trifolium repens}; plant diversity depended on topography and spatial effects rather than drainage, total N and C generally decreased with drainage and depth (but only interactively), the C/N ratio was drainage-independent, and $\delta^{13}$C depended significantly on drainage (depleted in drained plots) whereas $\delta^{15}$N did not \parencite{Harris2018}. In restored Iowa oxbows, tile drainage had negligible effects: physical variables did not differ significantly between multipurpose (tiled) and non-tiled oxbows, no biotic variable differed, and macrophyte, macroinvertebrate, and fish assemblages were similar (44 invertebrate taxa, 35 fish species, and 2{,}682 Topeka shiners recorded), providing preliminary but promising evidence that multipurpose oxbows can enhance biodiversity without ecological penalty \parencite{Leberg2024}. Wetland restoration outcomes improved with sediment excavation, which initially reduced invasive cover (notably hybrid cattail and reed canary grass) and increased diversity and evenness, though the invasive-suppression benefit was lost by $\sim$6 years; community composition differed significantly by wetland age, with younger and excavated wetlands favouring native perennials and submerged macrophytes \parencite{Winikoff2020}.

\subsection{Amendments, Cover Management, and Nutrient Retention: Additional Practices}

Two further practice classes complete the picture. Polyacrylamide (PAM) reduced downward phosphorus transport in soil columns, with the 0.20\% treatment optimal (cumulative soluble-P loss 348 versus 541\,mg\,tube$^{-1}$; mean soluble-P concentration 87 versus 138\,$\mu$g\,ml$^{-1}$ in the control); in field plots all PAM treatments significantly cut runoff volume and total-P loss and significantly reduced leachate volume to tile drains, though total-P losses in leachate were not significantly different (attributed to low surface-application rate and high variability); PAM alone maximised wheat growth and plant-P uptake, whereas PAM plus calcium amendments best reduced runoff-P losses \parencite{Jiang2010}. In Spanish citrus, combining anti-weed mesh with a drainage system improved vegetative growth, flower production, and yield relative to the traditional control: netted treatments generally produced larger canopies (the largest canopy volume 2256.17\,L and trunk diameter 40.38\,mm under one treatment) and higher juice yield, while the uncovered control produced the heaviest individual fruits; metabolomic profiling showed a shift from 27 to 25 detected metabolites between seasons (proline the most abundant amino acid, quinic acid the most abundant organic acid), with certain netted-and-drained treatments best maintaining metabolomic homeostasis, and the authors recommend anti-weed mesh for commercial plantations in water-scarce regions for its reduced water, herbicide, and labour costs \parencite{GiménezValero2024,GiménezValero2025}.